\documentclass[10pt,a4paper]{scrartcl}
\usepackage[T1]{fontenc}
\usepackage{newpxtext,newpxmath}
\usepackage[margin=12mm]{geometry}
\usepackage{multicol}
\usepackage{xcolor}
\definecolor{primary}{HTML}{1F4E79}
\setkomafont{section}{\normalsize\bfseries\color{primary}}
\RedeclareSectionCommand[beforeskip=7pt,afterskip=3pt]{section}
\setlength{\parindent}{0pt}
\setlength{\columnsep}{8mm}
\setlength{\columnseprule}{0.2pt}
\AtBeginDocument{%
\setlength{\abovedisplayskip}{4pt}\setlength{\belowdisplayskip}{4pt}%
\setlength{\abovedisplayshortskip}{3pt}\setlength{\belowdisplayshortskip}{3pt}}
\newcommand{\dd}{\mathrm{d}}
\newcommand{\ee}{\mathrm{e}}
\pagestyle{empty}
\raggedbottom
\begin{document}
\begin{center}
{\large\bfseries Physical Models of Biological Systems}\quad 77972\\[-1pt]
{\small Formula and concept sheet}
\end{center}
\vspace{-5pt}
\begin{multicols}{2}
\raggedcolumns
\small
\section*{Model formulation and inference}
\textbf{Cycle:} formulate $\to$ solve $\to$ validate $\to$ revise.\\
\textbf{Uses:} consistency, prediction, inference, intuition.\\
\textbf{Checks:} units, signs, initial values, positivity, limits.\\
\textbf{Steady state:} constant amounts; fluxes may persist.\\
\textbf{Identifiable:} distinct parameters give distinct data.\\
\textbf{Over-constrained:} more independent constraints than free parameters. Does not ensure identifiability.\\
\textbf{Overfit:} excess flexibility fits noise. Count independent features; correlated points add less information.

\section*{Production, clearance, and stability}
% Balance and exact relaxation for a constant source and linear removal.
\begin{gather*}
\dot W=Q_{\mathrm{in}}-Q_{\mathrm{out}},\qquad
\dot W=0\Rightarrow Q_{\mathrm{in}}=Q_{\mathrm{out}}\\
\dot W=Q-\kappa W,\qquad W^*=Q/\kappa\\
W(t)=W^*+[W(0)-W^*]\ee^{-\kappa t}
\end{gather*}
$[Q]=[W]/\text{time}$, $[\kappa]=\text{time}^{-1}$.\\
\textbf{Stable for $\kappa>0$:} deviations decay; $\tau=1/\kappa$.\\
\textbf{Steady amount fixes a ratio, not either flux.}
% Hydrostatic column and viscous outlet analogy.
\begin{equation*}
\Delta p=\rho gh,\quad W=Ah,\quad Q_{\mathrm{out}}\propto\Delta p\propto W
\end{equation*}
Fixed area $A$; density $\rho$; viscous Poiseuille outlet.

\section*{First-order decay and semilog plots}
% Constant hazard, characteristic times, and apparent decay rate.
\begin{gather*}
\dot N=-kN\Rightarrow N=N_0\ee^{-kt}\\
\tau=1/k,\qquad t_{1/2}=\ln2/k\\
\ln N=\ln N_0-kt,\qquad
\frac{\dd\log_{10}N}{\dd t}=-\frac{k}{\ln10}\\
r_{\mathrm{app}}=\frac{\ln[N(t_1)/N(t_2)]}{t_2-t_1}
\end{gather*}
For multiple modes, $r_{\mathrm{app}}$ need not be a process rate.

\section*{Linear forced equation: master method}
% Integrating factor normalized to one at the initial time.
\begin{gather*}
\dot y+a(t)y=f(t),\quad \mu(t)=\ee^{\int_0^t a(s)\dd s}\\
(\mu y)'=\mu f,\quad
y(t)=\frac{y(0)+\int_0^t\mu(s)f(s)\dd s}{\mu(t)}
\end{gather*}

\section*{Two-stage viral turnover}
\textbf{Assume:} no new infections after $t=0$; constant rates; mixed compartment; expected populations.\\
$N_I$: infected cells; $N_V$: total free particles, including noninfectious ones; $N_{I0},N_{V0}$: initial values.\\
$k_I,k_V$: clearance rates [time$^{-1}$]; $\gamma$: virions/(cell time); $\beta=\gamma N_{I0}$: virions/time. RNA load measures abundance, not infectivity.
% Source population and forced particle dynamics.
\begin{gather*}
\dot N_I=-k_IN_I,\qquad N_I=N_{I0}\ee^{-k_It}\\
\dot N_V=\gamma N_I-k_VN_V=\beta\ee^{-k_It}-k_VN_V
\end{gather*}
% General solution for unequal rates.
For $k_V\ne k_I$:
\begin{equation*}
N_V=N_{V0}\ee^{-k_Vt}
+\frac{\beta(\ee^{-k_It}-\ee^{-k_Vt})}{k_V-k_I}
\end{equation*}
% Initially balanced production and removal.
\textbf{Pretreatment balance:} $\beta=k_VN_{V0}$.
\begin{equation*}
\frac{N_V}{N_{V0}}=\frac{k_V\ee^{-k_It}-k_I\ee^{-k_Vt}}{k_V-k_I}
\end{equation*}
% Initial derivatives and local shoulder.
\begin{gather*}
\dot N_V(0)=0,\qquad \ddot N_V(0)=-k_Ik_VN_{V0}\\
N_V/N_{V0}=1-\tfrac12k_Ik_Vt^2+O(t^3)
\end{gather*}
\textbf{Shoulder:} horizontal tangent, not a finite flat interval.
% Equal-rate solution and initially steady special case.
For $k_I=k_V=k$:
\begin{gather*}
N_V=(N_{V0}+\beta t)\ee^{-kt}\\
\beta=kN_{V0}\Rightarrow N_V=N_{V0}(1+kt)\ee^{-kt}
\end{gather*}
\textbf{Equal rates still give zero initial slope.}
% Scale-separated tails.
\begin{align*}
k_V\gg k_I &: N_V\simeq(\beta/k_V)\ee^{-k_It}\\
k_I\gg k_V &: N_V\propto\ee^{-k_Vt}\\
t\to\infty &: -\dd\ln N_V/\dd t\to\min(k_I,k_V)
\end{align*}
\textbf{Inference:} measured $N_{V0}$ leaves 3 unknowns $(\beta,k_I,k_V)$; balance reduces to 2.\\
\textbf{Symmetry:} $k_I\leftrightarrow k_V$ preserves the balanced curve: only the unordered pair is identifiable. Also $\beta=\gamma N_{I0}$ cannot separate $\gamma,N_{I0}$.\\
\textbf{Scope:} early treatment response; omits reservoirs and long-term progression.


\section*{Pulse-chase: source and destination}
% A labelled cohort moves while every compartment total stays steady.
\begin{gather*}
\dot D=kS-dD,\qquad D=K*S\\
D=k\!\int_{t_0}^{t}\!\ee^{-d(t-t')}S(t')\dd t'
  +\ee^{-d(t-t_0)}D_0\\
K(\tau)=k\ee^{-d\tau},\qquad \tilde D/\tilde S=k/(s+d)
\end{gather*}
$S,D$: label per cell upstream and downstream; $k$: transfer rate; $d$: labelled loss rate, both [time$^{-1}$].\\
\textbf{Fit:} with $G_d(t_i)=\int_{t_0}^{t_i}\ee^{-d(t_i-t')}S\dd t'$,
\begin{equation*}
k_*=\frac{\sum_iD_iG_i}{\sum_iG_i^2},\qquad
\chi^2_{\min}=\sum_iD_i^2-\frac{\bigl(\sum_iD_iG_i\bigr)^2}{\sum_iG_i^2}
\end{equation*}
\textbf{Fit the integral, not the slope:} differentiating multiplies noise by $\mathrm{i}\omega$, integrating divides by it.\\
\textbf{Division dilutes:} add $\ln2\,r_{\mathrm{div}}$ to the effective loss rate.

\section*{HIV: infectious and non-infectious virions}
% Infection is the only bilinear term; protease inhibition closes the infectious branch.
\begin{align*}
\dot N_U&=\lambda-k_UN_U-\epsilon N_VN_U\\
\dot N_I&=\epsilon N_VN_U-k_IN_I\\
\dot N_V&=\epsilon'\gamma N_I-k_VN_V\\
\dot N_X&=(1-\epsilon')\gamma N_I-k_VN_X
\end{align*}
$\epsilon$: infection coefficient, [(virion$\cdot$time)$^{-1}$]; $\epsilon'$: infectious fraction, dimensionless.
\begin{gather*}
N_U^*=\frac{k_Ik_V}{\epsilon\epsilon'\gamma},\qquad
\mathcal R_0=\frac{\lambda\epsilon\epsilon'\gamma}{k_Uk_Ik_V}
\end{gather*}
Infected state exists only for $\mathcal R_0>1$. \textbf{Protease inhibitor:} $\epsilon'\to0$ gives $N_V=N_{V0}\ee^{-k_Vt}$.
\section*{Independent resistance barriers}
% Simultaneous independent mutations in one replication opportunity.
\begin{equation*}
P(\text{all }n)=p^n,\quad
\mathbb E[\text{resistant copies}]=Mp^n
\end{equation*}
$p$: each mutation probability; $M$: opportunities. Requires independence and simultaneous mutations; omits sequential evolution and cross-resistance.

\section*{Discrete probability}
\begin{gather*}
\sum_\ell P(\ell)=1,\quad
P(E_1\cup E_2)=P(E_1)+P(E_2)-P(E_1\cap E_2)\\
P(E\mid E')=\frac{P(E\cap E')}{P(E')},\quad
P(E\cap E')=P(E)P(E')\ \text{(independent)}\\
P(H\mid E)=\frac{P(E\mid H)P(H)}{\sum_jP(E\mid H_j)P(H_j)}
\end{gather*}
\textbf{Geometric:} $P(J=j)=\xi(1-\xi)^{j-1}$, $\mathbb E[J]=1/\xi$, $\mathrm{Var}=(1-\xi)/\xi^2$.\\
\textbf{Binomial:} $P(K=k)=\binom nk\xi^k(1-\xi)^{n-k}$, $\mathbb E[K]=n\xi$.\\
\textbf{Screening:} prevalence $\pi$, sensitivity $s$, specificity $c$,
\begin{equation*}
P(D\mid+)=\frac{\pi s}{\pi s+(1-\pi)(1-c)}
\end{equation*}
Rare disease: a positive test is still mostly wrong.

\section*{Entropy and mutual information}
\begin{gather*}
H(X)=-\sum_xP(x)\log_2P(x),\qquad 0\le H\le\log_2M\\
H(Y\mid X)=\sum_xP(x)\,H(Y\mid X=x)\\
I(X;Y)=H(X)+H(Y)-H(X,Y)=H(Y)-H(Y\mid X)\\
I(X;Y)=\sum_{x,y}P(x,y)\log_2\frac{P(x,y)}{P(x)P(y)}\ \ge0\\
D_{\mathrm{KL}}(P\Vert Q)=\sum_xP(x)\log_2\frac{P(x)}{Q(x)}\ \ge0
\end{gather*}
$H$ in bits; $2^H$ is an effective number of states. $I=0$ iff independent; $D_{\mathrm{KL}}=0$ iff $P=Q$, and it is asymmetric.\\
\textbf{Jensen:} $\phi(\mathbb E Z)\le\mathbb E[\phi(Z)]$ for convex $\phi$; taking $\phi=-\log_2$ proves both inequalities.\\
\textbf{Bernoulli KL:} $D(p\Vert q)=p\log_2\frac pq+(1-p)\log_2\frac{1-p}{1-q}$.\\
\textbf{BSC, error $\varepsilon$:} $I=H(Y)-h(\varepsilon)\le1-h(\varepsilon)$, $h(\varepsilon)=-\varepsilon\log_2\varepsilon-(1-\varepsilon)\log_2(1-\varepsilon)$.\\
\textbf{Data processing:} $X\to Y\to Z$ implies $I(X;Z)\le I(X;Y)$.
\end{multicols}
\end{document}
